\section{Conclusion, limitations, future work and viva preparation}

\subsection{Limitations (state these honestly)}
\begin{itemize}
  \item \textbf{Ablation ties on real data.} On Intel Lab and PEMS-BAY, the model without message passing performs as
        well as the full GNN. The real-data gain comes mainly from the learned residual decoder (the correction on top of
        Kalman-greedy and the pairwise term); message passing helps measurably only on synthetic layouts.
  \item \textbf{Large networks.} On synthetic networks with 100 or more sensors (beyond the training range), Kalman-greedy
        is 8--12\% better than the GNN. Training on larger networks is the obvious fix.
  \item \textbf{Smaller gains on traffic.} PEMS-BAY sensors are less redundant than lab thermometers, so the gain is
        2.4\% in RMSE.
  \item \textbf{Model assumptions.} The sink assumes linear-Gaussian AR(1) dynamics; the channel is ideal apart from
        capacity and interference (no packet loss).
  \item \textbf{Offline training.} The correlation model is fitted once; it does not adapt if the environment changes.
\end{itemize}

\begin{saybox}[If asked: ``What does not work?'']
``Two things. On real data, removing the message passing does not hurt, so the improvement there comes from the
learned decoder rather than the graph layers. And on simulated networks larger than we trained on, plain Kalman-greedy is
still about 10 percent better. Both point to clear next steps.''
\end{saybox}

\subsection{Conclusion}
\begin{itemize}
  \item We identified, from the literature, that correlation-aware AoI work is analytical and hand-modelled, learned
        schedulers use interference graphs and soft freshness limits, and real-data evidence is scarce.
  \item We built \textbf{GNN-AoI}: a scheduler that learns the correlation structure from data, tracks the field with a
        Kalman sink, chooses transmissions with a GNN pointer decoder initialised as Kalman-greedy and trained with PPO,
        and guarantees every sensor's freshness by construction.
  \item On \textbf{real held-out data} it lowers the sink's error by \textbf{20\%} (Intel Lab, 45 sensors) and
        \textbf{2.4\%} (PEMS-BAY, 325 sensors) compared with AoI-based scheduling, and beats Kalman-greedy in
        \textbf{every} test window of both datasets.
  \item On simulated fields it halves the error of AoI scheduling and beats greedy selection that knows the true
        covariance.
  \item It runs in milliseconds, comes as a tested open-source package with reproducible experiments, and includes a
        live demo of real transmissions.
  \item The central lesson: \textbf{fresher is not the same as better-informed}. Using correlations to decide who
        transmits gives the sink a much better picture with the same channel, as long as freshness is still
        guaranteed.
\end{itemize}

\subsection{Future work}
\begin{itemize}
  \item Lossy channels with retransmissions (HARQ) \cite{6}.
  \item Online re-fitting of the correlation graph as seasons change or sensors are added.
  \item Training on larger networks; sparse message passing for thousands of sensors.
  \item Non-linear sinks (graph-recurrent forecasters such as DCRNN \cite{24}) in place of the Kalman filter.
  \item Distributed scheduling: each gateway runs the same shared GNN locally.
  \item A hardware pilot on a LoRa or Zigbee testbed.
\end{itemize}

\subsection{Likely viva questions}
\begin{description}[style=nextline,leftmargin=0.6cm]
  \item[What is Age of Information?] Slots since the sink last heard from a sensor. It measures freshness, not how much
        the sink knows.
  \item[Why not just minimise AoI?] Because correlated sensors make some transmissions redundant. Max-AoI has the
        lowest age but twice the error of our method (Figure ``fresher is not better-informed'').
  \item[What exactly does the GNN output?] A probability for each sensor at each pick step; the top choice is taken,
        the Kalman covariance is updated, and the next sensor is chosen, until the slot is full.
  \item[What is the reward?] $-(1-\gamma)\log(\text{error}+0.005)$ per slot: lower error, higher reward.
  \item[Why PPO and not DQN?] Choosing 6 of 45 sensors has millions of combinations; DQN needs a value per action. PPO
        trains our step-by-step selection policy directly and stably.
  \item[Why does the GNN start as Kalman-greedy?] So it never starts worse than the strongest classical method; RL only
        learns improvements. Behaviour cloning, which we tried first, never matched the teacher.
  \item[How do you guarantee freshness?] Sensors at age $F$ are forced into the slot before the scheduler acts; the
        environment asserts it every slot, and tests check every scheduler.
  \item[How do you avoid data leakage?] Chronological split; the profile, standardisation and model come only from the
        training period; validation uses separate data; a unit test checks it.
  \item[Why two datasets?] One with strong spatial correlation (indoor temperature), one large and less redundant
        (freeway traffic): they show where the method helps a lot and where it helps a little.
  \item[Why does Kalman-greedy lose on PEMS-BAY?] Its fitted model is over-confident: it believes it knows sensors it
        does not, so it polls them too rarely. Our policy saw this mismatch during training on real replays.
  \item[Is the improvement statistically significant?] Yes: Intel 10/10 windows vs.\ Kalman-greedy ($p=0.002$), PEMS-BAY
        15/15 ($p=6\times10^{-5}$), synthetic 73\% of layouts ($p=0.008$), Wilcoxon signed-rank tests.
  \item[How long does a decision take?] 4.6\,ms (Intel), 37\,ms (PEMS-BAY) on a laptop CPU; the slot is 5 minutes.
  \item[What is your base paper and what did you add?] Ramakanth, Tripathi \& Modiano, INFOCOM 2024. We add multiple polls
        per slot, correlations learned from data, a hard freshness bound, a learned GNN policy and real-data evaluation.
\end{description}

\subsection{References}\label{sec:refs}
{\small
\begin{enumerate}[label={[\arabic*]},leftmargin=1cm,itemsep=0.15em]
  \item V.~Tripathi and E.~Modiano, ``Optimizing age of information with correlated sources,'' \emph{Proc.\ ACM MobiHoc}, 2022, pp.~41--50; extended in \emph{IEEE/ACM Trans.\ Netw.}, vol.~32, no.~6, pp.~4660--4675, 2024.
  \item R.~V.~Ramakanth, V.~Tripathi, and E.~Modiano, ``Monitoring correlated sources: AoI-based scheduling is nearly optimal,'' \emph{Proc.\ IEEE INFOCOM}, 2024, pp.~641--650; extended in \emph{IEEE Trans.\ Mobile Comput.}, vol.~24, no.~2, pp.~1043--1054, 2025.
  \item E.~Erbayat, A.~Maatouk, P.~Zou, and S.~Subramaniam, ``Age of information optimization and state error analysis for correlated multi-process multi-sensor systems,'' \emph{Proc.\ ACM MobiHoc}, 2024, pp.~331--340.
  \item E.~Erbayat, A.~Maatouk, P.~Zou, and S.~Subramaniam, ``Age of information optimization with preemption strategies for correlated systems,'' \emph{Proc.\ IEEE ISIT}, 2025.
  \item M.~Fidler, F.~Gallistl, J.~P.~Champati, and J.~Widmer, ``2D-AoI: Age-of-information of distributed sensors for spatio-temporal processes,'' \emph{IEEE Trans.\ Commun.}, vol.~74, pp.~645--661, 2026.
  \item E.~T.~Ceran, D.~G\"und\"uz, and A.~Gy\"orgy, ``A reinforcement learning approach to age of information in multi-user networks with HARQ,'' \emph{IEEE J.\ Sel.\ Areas Commun.}, vol.~39, no.~5, pp.~1412--1426, 2021.
  \item H.~Farag, M.~Gidlund, and \v{C}.~Stefanovi\'c, ``A deep reinforcement learning approach for improving age of information in mission-critical IoT,'' \emph{Proc.\ IEEE GCAIoT}, 2021, pp.~14--18.
  \item J.~Chen, W.~Liu, D.~E.~Quevedo, Y.~Li, and B.~Vucetic, ``Structure-enhanced deep reinforcement learning for optimal transmission scheduling,'' \emph{Proc.\ IEEE ICC}, 2023, pp.~1212--1218.
  \item B.~G.~Reyhan and S.~Coleri, ``Safe deep reinforcement learning for resource allocation with peak age of information violation guarantees,'' \emph{IEEE Trans.\ Commun.}, vol.~73, no.~12, pp.~14197--14211, 2025.
  \item Y.~Shen, Y.~Shi, J.~Zhang, and K.~B.~Letaief, ``Graph neural networks for scalable radio resource management: Architecture design and theoretical analysis,'' \emph{IEEE J.\ Sel.\ Areas Commun.}, vol.~39, no.~1, pp.~101--115, 2021.
  \item Z.~Zhao, G.~Verma, C.~Rao, A.~Swami, and S.~Segarra, ``Link scheduling using graph neural networks,'' \emph{IEEE Trans.\ Wireless Commun.}, vol.~22, no.~6, pp.~3997--4012, 2023.
  \item L.~Wang, Q.~Wang, H.~Chen, and S.~Zhou, ``Age of information-oriented probabilistic link scheduling for device-to-device networks,'' \emph{Proc.\ IEEE/IFIP WiOpt}, 2024 (arXiv:2410.20196).
  \item W.~Wang, Q.~Wu, P.~Fan, N.~Cheng, W.~Chen, J.~Wang, and K.~B.~Letaief, ``Optimizing age of information in vehicular edge computing with federated graph neural network multi-agent reinforcement learning,'' \emph{IEEE Trans.\ Wireless Commun.}, vol.~25, pp.~23017--23032, 2026.
  \item R.~D.~Yates, Y.~Sun, D.~R.~Brown, S.~K.~Kaul, E.~Modiano, and S.~Ulukus, ``Age of information: An introduction and survey,'' \emph{IEEE J.\ Sel.\ Areas Commun.}, vol.~39, no.~5, pp.~1183--1210, 2021.
  \item S.~Kaul, R.~Yates, and M.~Gruteser, ``Real-time status: How often should one update?'' \emph{Proc.\ IEEE INFOCOM}, 2012, pp.~2731--2735.
  \item I.~Kadota, A.~Sinha, E.~Uysal-Biyikoglu, R.~Singh, and E.~Modiano, ``Scheduling policies for minimizing age of information in broadcast wireless networks,'' \emph{IEEE/ACM Trans.\ Netw.}, vol.~26, no.~6, pp.~2637--2650, 2018.
  \item A.~Maatouk, S.~Kriouile, M.~Assaad, and A.~Ephremides, ``The age of incorrect information: A new performance metric for status updates,'' \emph{IEEE/ACM Trans.\ Netw.}, vol.~28, no.~5, pp.~2215--2228, 2020.
  \item A.~S.~Leong, A.~Ramaswamy, D.~E.~Quevedo, H.~Karl, and L.~Shi, ``Deep reinforcement learning for wireless sensor scheduling in cyber--physical systems,'' \emph{Automatica}, vol.~113, art.~108759, 2020.
  \item A.~Krause, A.~Singh, and C.~Guestrin, ``Near-optimal sensor placements in Gaussian processes: Theory, efficient algorithms and empirical studies,'' \emph{J.\ Mach.\ Learn.\ Res.}, vol.~9, pp.~235--284, 2008.
  \item R.~E.~Kalman, ``A new approach to linear filtering and prediction problems,'' \emph{Trans.\ ASME J.\ Basic Eng.}, vol.~82, no.~1, pp.~35--45, 1960.
  \item T.~N.~Kipf and M.~Welling, ``Semi-supervised classification with graph convolutional networks,'' \emph{Proc.\ ICLR}, 2017.
  \item W.~Kool, H.~van Hoof, and M.~Welling, ``Attention, learn to solve routing problems!'' \emph{Proc.\ ICLR}, 2019.
  \item J.~Schulman, F.~Wolski, P.~Dhariwal, A.~Radford, and O.~Klimov, ``Proximal policy optimization algorithms,'' arXiv:1707.06347, 2017.
  \item Y.~Li, R.~Yu, C.~Shahabi, and Y.~Liu, ``Diffusion convolutional recurrent neural network: Data-driven traffic forecasting,'' \emph{Proc.\ ICLR}, 2018 (PEMS-BAY dataset).
  \item P.~Bodik, W.~Hong, C.~Guestrin, S.~Madden, M.~Paskin, and R.~Thibaux, ``Intel Lab Data,'' 2004. Online: \url{http://db.csail.mit.edu/labdata/labdata.html}
\end{enumerate}
}
